\section{Admissible extensions with prescribed support}
Let a finite group $A$ act on finite sets $Z$ and $E$, and let
$p:Z\to E$ be equivariant. Let $R\subseteq Z$ be invariant; it is the set
passing the application's admissibility tests. Suppose a finite poset $L$
with greatest element $1$ has an $A$-action by order automorphisms and an
equivariant map $s:R\to L$. Write $R_l=\{z\in R:s(z)\le l\}$.
The individual sets $R_l$ need not be invariant. Choose one $e$ in each
$A$-orbit in $p(R)$, and put $R(e)=R\cap p^{-1}(e)$.

\textbf{Theorem UD.AdmissibleExtensions.}
The equivalence classes of admissible extensions of exact support $1$ are
in bijection with the disjoint union of the $A_e$-orbits on
$\{z\in R(e):s(z)=1\}$. Their number is
\[
 \sum_{e\in A\backslash p(R)}\frac1{|A_e|}
 \sum_{a\in A_e}\sum_{l\in L}\mu(l,1)
       |\Fix_{R(e)}(a)\cap R_l|,
\]
where $\mu$ is the M\"obius function of $L$.

\emph{Proof.}
The greatest element is fixed by every order automorphism, so the exact
support condition is invariant. Apply UD.ExtensionOrbits to the restriction
of $p$ to that subset. Its fiber equivalence is exactly the $A_e$-action:
an element carrying one extension over $e$ to another necessarily fixes $e$.
An empty exact-support fiber contributes zero, so representatives may be
taken from all of $p(R)$.

For each $z\in R$, the defining inverse identity of the incidence-algebra
M\"obius function gives
\[
 \sum_{s(z)\le l\le1}\mu(l,1)=
 \begin{cases}1,&s(z)=1,\\0,&s(z)<1.\end{cases}
\]
Sum over the points of $R(e)$ fixed by a particular $a\in A_e$. This gives
the inner sum in the formula. Finally, for a finite group $H$ acting on a
finite set $Y$, count the pairs $(h,y)$ with $hy=y$. Their number is both
$\sum_h|\Fix_Y(h)|$ and $\sum_y|H_y|$. Each orbit contributes $|H|$ to the
second sum by the orbit--stabilizer identity. Division by $|H|$ gives the
number of orbits. Use this identity with $H=A_e$ and the exact-support
fiber, and sum over the representatives $e$. \(\square\)

\textbf{Corollary UD.ExactSupportOrbits.}
If $A$ acts on a finite set $X$ with an equivariant support map to $L$,
the number of exact-top-support orbits is
\[
 \frac1{|A|}\sum_{a\in A}\sum_{l\in L}
       \mu(l,1)|\Fix_X(a)\cap\{x:s(x)\le l\}|.
\]
\emph{Proof.} Take $E$ to be a one-point set and $R=Z=X$.
\(\square\)

Taking a one-element support poset instead recovers counting in ordinary
extension fibers. For an involution $D$ on a finite set $Y$, it gives
$(|Y|+|\Fix_Y(D)|)/2$ orbits. A division by two is therefore justified
without a correction term only after freeness has been proved.

\textbf{Lemma UD.OddOrbits.}
Let a finite group $A$ of odd order act on $X$. Let $D$ be a fixed-point-free
involution of $X$ commuting with this action. Then $D$ has no fixed orbit
in $A\backslash X$.

\emph{Proof.} If $Dx=ax$, commutation and $D^2=1$ give $a^2x=x$.
The order of $a$ is odd, so $a$ is a power of $a^2$ and fixes $x$.
Then $Dx=x$, a contradiction. \(\square\)

For the divisor poset of an integer $N$, support often means the least
field $\mathbb F_{2^d}$ containing all coordinates, where $d\mid N$.
The application must establish that its subfield sets are exactly the
support sublevel sets. If $A(d)$ counts data defined over that field,
the number of data of exact degree $N$ is
\[
 E(N)=\sum_{d\mid N}\mu(N/d)A(d).
\]
If Frobenius acts freely on them with orbit length $N$, division by $N$
counts its orbits. For odd $N$, a commuting fixed-point-free involution
then permits another division by two, by UD.OddOrbits. These are conditional
consequences of the common theorem, not substitutes for the generation and
stabilizer calculations in the applications.
